\section{Section balance and rigidity}\label{sec:rigidity}

We now convert the exclusion of models with two or more directions into
uniform control of tangent sections.  This control supplies the modulus
of convexity needed for the interior estimates of Trudinger and Wang.

\begin{lemma}[Uniform balance of tangent sections]\label{lem:uniform-balance}
There is a constant \(0<\rho\le1\) such that, for every
\(a\in\Omega\) and \(t>0\), the compact convex body
\[
 K_a(t)=\overline{S_u(a,t)}-a
\]
satisfies
\begin{equation}\label{eq:uniform-section-balance}
 -\rho K_a(t)\subset K_a(t).
\end{equation}
The same \(\rho\) works for all base points and heights.
\end{lemma}

\begin{proof}
For a fixed \(a\), use the ambient affine coordinates
\[
 s=z-l_a(x),\qquad y=x-a.
\]
The epigraph lies in \(s\ge0\), contains the ray
\(\{s\ge0,\ y=0\}\), and has compact caps by
Lemma~\ref{lem:proper-sections}.  It is therefore a model with one
direction, whose fiber at height \(t\) is \(K_a(t)\).
Proposition~\ref{prop:one-direction} says that the maximal direction
number is one.  The uniformity over models and coordinates in
Proposition~\ref{prop:centering} consequently gives
\(-K_a(t)\subset C K_a(t)\) with one finite constant \(C\).
Taking \(\rho=\min\{1,C^{-1}\}\) proves the claim.
\end{proof}

\begin{lemma}[Doubling and extension of rays]\label{lem:doubling-entire}
The domain is \(\Omega=\R^n\).  Moreover, there are constants
\(\beta>1\) and \(A>1\), depending only on \(\rho\), such that every
tangent-subtracted line restriction
\[
 g(r)=u(a+re)-u(a)-r\,Du(a)\cdot e
\]
satisfies
\begin{equation}\label{eq:gap-doubling}
 g(\beta r)\le A g(r)\qquad(r>0).
\end{equation}
One may take
\begin{equation}\label{eq:doubling-constants}
 \beta=1+\frac{\rho}{2},\qquad
 A=\frac{(2+\rho)(1+\rho)}{\rho}.
\end{equation}
\end{lemma}

\begin{proof}
Initially let \(I=\{r:a+re\in\Omega\}\), an open interval containing
zero, and choose \(R\in I\) with \(R>0\).
For \(0<r_0<R\), the tangent section based at \(a+r_0e\) and having
height
\[
 b=r_0g'(r_0)-g(r_0)>0
\]
contains \(a\) on its boundary.  Applying
\eqref{eq:uniform-section-balance} to the displacement \(-r_0e\)
puts the displacement \(\rho r_0e\) in the same translated section.
In particular \((1+\rho)r_0\in I\), and the section inequality gives
\[
 g((1+\rho)r_0)\le(1+\rho)r_0g'(r_0).
\]
Convexity and \(g\ge0\) give
\[
 g'(r_0)\le
 \frac{g(R)-g(r_0)}{R-r_0}
 \le\frac{g(R)}{R-r_0}.
\]
Set
\[
 \theta=\frac{2+\rho}{2(1+\rho)},\qquad r_0=\theta R.
\]
Then \((1+\rho)\theta=\beta\), with \(\beta\) as in
\eqref{eq:doubling-constants}, and
\[
 \beta R\in I,\qquad
 g(\beta R)\le\frac{\beta}{1-\theta}g(R)=A g(R).
\]
Iterating the domain inclusion shows that the positive ray from \(a\)
in direction \(e\) is contained in \(\Omega\).
Every oriented direction admits an initial positive interval by
openness, so all rays from \(a\) lie in \(\Omega\).  This proves
\(\Omega=\R^n\), as well as \eqref{eq:gap-doubling} on every ray.
\end{proof}

Fix a base point, translate it to zero, and subtract its tangent, so that
\(u(0)=0\), \(Du(0)=0\), and \(u\ge0\).  Each section \(S_u(0,t)\)
can be normalized by a linear map that preserves zero, with constants
uniform in \(t\).  For \(K=\overline{S_u(0,t)}\), apply
Lemma~\ref{lem:linear-normalization} with \(d=n\) and
\(C_*=\rho^{-1}\).  It gives an invertible linear map \(L\) with
\[
 \overline{B_{2\rho/(1+\rho)}}\subset LK
 \subset\overline{B_{2n}}.
\]
A scalar adjustment therefore produces invertible maps \(A_t\) and
a fixed \(R\ge1\), depending only on \(n,\rho\), for which
\begin{equation}\label{eq:normalized-sections}
 v_t(y)=t^{-1}u(A_ty),\qquad
 B_1\subset S_{v_t}(0,1)\subset B_R.
\end{equation}
The scalar adjustment may be made with a strict margin, so the inner
inclusion holds for the open section.  Balance and
\eqref{eq:gap-doubling} are unchanged by these horizontal linear
maps and positive vertical rescalings.

\begin{lemma}[A uniform normalized modulus]\label{lem:normalized-modulus}
For all functions \(v_t\) in \eqref{eq:normalized-sections}, there is one
positive function \(b:(0,\infty)\to(0,\infty)\) such that
\begin{equation}\label{eq:normalized-modulus}
 S_{v_t}(x,b(r))\subset x+B_r
 \qquad\bigl(x\in B_{1/4},\ r>0\bigr).
\end{equation}
In particular, with \(\gamma=(4R)^{-1}\),
\begin{equation}\label{eq:uniform-strict-convexity}
 S_u(x,tb(r))\subset x+rS_u(0,t)
 \qquad
 \bigl(x\in\gamma S_u(0,t),\ r,t>0\bigr).
\end{equation}
\end{lemma}

\begin{proof}
Write \(v=v_t\).  Since \(v(0)=Dv(0)=0\) and \(v\le1\) on
\(\overline{B_1}\), convexity gives \(0\le v(x)\le1/4\) on
\(B_{1/4}\).  Fix such an \(x\).
For every unoriented line through \(x\), choose its unit direction
\(e\) so that \(Dv(x)\cdot e\le0\).
The ray in this direction meets \(\partial S_v(0,1)\) at a distance
at most \(2R\).  At that intersection, its tangent-subtracted gap
\[
 g_{x,e}(r)=v(x+re)-v(x)-r\,Dv(x)\cdot e
\]
is at least \(1-v(x)\ge3/4\).
Thus the positive-ray radius of the closed tangent section of height
\(1/4\), based at \(x\), is at most \(2R\).
Balance of that section bounds the opposite-ray radius by
\(2R/\rho\): a point on the opposite ray, reflected and contracted
by \(\rho\), must still fit on the first ray.
Consequently every unit direction \(e\) satisfies
\begin{equation}\label{eq:fixed-radius-gap}
 g_{x,e}(R_*)\ge\frac14,\qquad R_*=\frac{2R}{\rho}.
\end{equation}

The gap \(g_{x,e}\) is nondecreasing on the positive ray.  Define
\begin{equation}\label{eq:explicit-modulus}
 N(r)=\max\left\{0,\left\lceil\log_\beta\frac{R_*}{r}\right\rceil\right\},
 \qquad b(r)=\frac{1}{8A^{N(r)}}.
\end{equation}
Iterating \eqref{eq:gap-doubling} and using
\eqref{eq:fixed-radius-gap} gives
\[
 A^{N(r)}g_{x,e}(r)
 \ge g_{x,e}\bigl(\beta^{N(r)}r\bigr)
 \ge\frac14.
\]
Thus the tangent gap is at least \(2b(r)\) at every point of
distance \(r\) from \(x\), and at every point farther away.
This proves \eqref{eq:normalized-modulus}.

If \(x\in\gamma S_v(0,1)\), then \(|x|<1/4\).  Since
\(B_1\subset S_v(0,1)\), we obtain
\[
 S_v(x,b(r))\subset x+B_r\subset x+rS_v(0,1).
\]
Applying \(A_t\) gives \eqref{eq:uniform-strict-convexity}.
All constants and the function \(b\) are independent of the height
\(t\), the base \(x\) in the stated range, and the line direction.
\end{proof}

Condition \eqref{eq:uniform-strict-convexity} is precisely the
uniform strict convexity condition of
\cite[Condition~(4.10)]{TW2000}, with base point zero.
For the final step we use its underlying interior estimate, so that
the effect of the section rescaling is explicit.

\begin{theorem}[Trudinger--Wang interior estimate]\label{thm:tw-interior}
Let \(O\subset\R^n\) be a bounded convex domain and let
\(f\in C^4(O)\) have positive-definite Hessian and solve the affine
maximal equation.  Suppose \(-1\le f\le0\).
If a positive function \(b_0\) gives the section control
\[
 S_{f|_O}(x,b_0(r))\subset B_r(x)
 \qquad(x\in O,\ r>0),
\]
then \(f\in C^\infty(O)\), and for each \(O'\Subset O\) and
integer \(q\ge2\) there are
positive constants \(c,C_q\) such that
\[
 D^2f\ge c\Id,\qquad |D^qf|\le C_q
 \quad\hbox{on }O'.
\]
The constants depend only on \(n\), \(\operatorname{diam}O\),
\(\operatorname{dist}(O',\partial O)\), and the specified lower
modulus \(b_0\), with additional dependence on \(q\) for \(C_q\).
\end{theorem}

This is \cite[Theorem 4.2]{TW2000}, stated in terms of a prescribed
lower bound for the modulus of convexity.  Its dimension is
arbitrary.  The sections \(S_{f|_O}\) here are taken within \(O\).
Its proof combines the linearized Monge--Amp\`ere estimates of
Caffarelli and Guti\'errez~\cite{CaffarelliGutierrez1997},
Caffarelli's Monge--Amp\`ere regularity theory~\cite{Caffarelli1990},
and Schauder estimates; see \cite[Section~4]{TW2000}.
The rescaling in the following proof is the reduction in
\cite[Theorem~2.1 and its proof]{TW2000}, with the uniform modulus now
supplied by Lemma~\ref{lem:normalized-modulus}.

\begin{proof}[Proof of Theorem~\ref{thm:main}]
Proposition~\ref{prop:one-direction} applies for every \(3\le n\le9\).
Lemmas~\ref{lem:uniform-balance} and \ref{lem:doubling-entire}
therefore give \(\Omega=\R^n\) and uniform section balance.
Fix any base point, translate it to zero, subtract its tangent, and
use the normalized functions \(v_t\) from
\eqref{eq:normalized-sections}.

These functions solve the same affine maximal equation.
Indeed, for \(v(y)=t^{-1}u(Ay)\), set
\[
 c_t=t^{-n}(\det A)^2,\qquad a=\frac{n+1}{n+2}.
\]
Then
\[
 w_v=c_t^{-a}w_u\circ A,\qquad
 U_v=t^{1-n}(\det A)^2 A^{-1}(U_u\circ A)A^{-\mathsf T},
\]
so the contraction \(U_v^{ij}(w_v)_{ij}\) is a nonzero constant
multiple of \((U_u^{ij}(w_u)_{ij})\circ A\), and vanishes.
Subtracting an affine function leaves the equation unchanged as well.

Apply Theorem~\ref{thm:tw-interior} to \(f_t=v_t-1\) on
\(O=B_{1/4}\), with \(O'=B_{1/8}\).
There \(-1\le f_t\le-3/4\), and
Lemma~\ref{lem:normalized-modulus} supplies the same lower modulus
at every point of \(O\).  Restricting a section to \(O\) only
decreases it.  We obtain constants independent of \(t\) such that
\begin{equation}\label{eq:normalized-derivative-bounds}
 D^2v_t(0)\ge c_0\Id,\qquad |D^3v_t(0)|\le C_0.
\end{equation}

Put \(H=D^2u(0)>0\).  The identity
\[
 D^2v_t(0)=t^{-1}A_t^{\mathsf T}HA_t
\]
and the first bound in \eqref{eq:normalized-derivative-bounds} imply
\[
 t c_0|z|^2
 \le (A_tz)^{\mathsf T}H(A_tz)
 \le\lambda_{\max}(H)|A_tz|^2,
\]
and hence
\begin{equation}\label{eq:inverse-normalization-bound}
 \|A_t^{-1}\|
 \le\sqrt{\lambda_{\max}(H)/c_0}\,t^{-1/2}.
\end{equation}
For any three vectors \(z_1,z_2,z_3\), differentiation of
\(u(x)=t\,v_t(A_t^{-1}x)\) gives
\[
 D^3u(0)[z_1,z_2,z_3]
 =t\,D^3v_t(0)
       [A_t^{-1}z_1,A_t^{-1}z_2,A_t^{-1}z_3].
\]
Equations \eqref{eq:normalized-derivative-bounds} and
\eqref{eq:inverse-normalization-bound} therefore give
\[
 |D^3u(0)|\le Ct^{-1/2}\longrightarrow0
 \qquad\text{as }t\longrightarrow\infty.
\]
The base point was arbitrary, so \(D^3u\equiv0\).  Thus
\[
 u(x)=\tfrac12 x^{\mathsf T}Qx+b\cdot x+c
\]
on \(\R^n\), with \(Q>0\).
Its graph is an invertible affine image of
\(\{(z,|z|^2):z\in\R^n\}\), as required.
\end{proof}

\subsection*{The affine-complete consequence}

The graph theorem now gives the consequence announced in the
introduction. The completeness and global graph results of Trudinger
and Wang represent the hypersurface as a complete convex graph.

\begin{corollary}[Affine-complete hypersurfaces]\label{cor:affine-complete}
Let \(3\le n\le9\). Let \(M\) be a smooth locally uniformly convex
immersed hypersurface in \(\R^{n+1}\) whose underlying manifold is
connected and open (noncompact and without boundary). Suppose that
\(M\) is complete for the affine (Berwald--Blaschke) metric and is
classically affine maximal (stationary for affine area under compactly
supported smooth variations). Then \(M\) is an invertible affine image
of the standard elliptic paraboloid.
\end{corollary}

\begin{proof}
By \cite[Theorem~A]{TW2002}, \(M\) is complete for the induced
Euclidean metric. By \cite[Corollary~1]{TW2002}, an invertible affine
change of coordinates gives an affine image \(\widetilde M\) that is
the global graph of a convex function \(u\) over a nonempty open convex
domain \(\Omega\subset\R^n\), not assumed to equal \(\R^n\).
The linear part \(L\) of this change uniformly compares Euclidean
lengths by its least and greatest singular values, preserving
completeness. By Equation~\eqref{eq:area-covariance}, it multiplies
affine area on compact patches by the fixed positive factor
\(|\det L|^{n/(n+2)}\), preserving stationarity.
Smoothness and local uniform convexity give
\(u\in C^\infty(\Omega)\) with positive-definite Hessian, and classical
stationarity gives \eqref{eq:affine-maximal}. Theorem~\ref{thm:main}
now implies \(\Omega=\R^n\) and that \(u\) is a quadratic polynomial
with positive-definite Hessian. Undoing the affine change of coordinates
proves the conclusion.
\end{proof}
